\documentclass[11pt,a4paper]{article}
\usepackage[utf8]{inputenc}
\usepackage[french]{babel}
\usepackage[T1]{fontenc}
\usepackage{geometry}
\geometry{margin=2.5cm}
\usepackage{xcolor}
\usepackage{hyperref}
\usepackage{enumitem}
\usepackage{fancyhdr}
\usepackage{tcolorbox}
\usepackage{titlesec}
\usepackage{charter} % Modern clean font

% Define professional color palette
\definecolor{primary}{HTML}{3B82F6}     % Bright Blue
\definecolor{secondary}{HTML}{1E3A8A}   % Deep Blue
\definecolor{dark}{HTML}{0F172A}        % Slate Dark
\definecolor{light}{HTML}{F8FAFC}       % Light Slate BG
\definecolor{border}{HTML}{E2E8F0}      % Light Slate Border

% Page style setup
\pagestyle{fancy}
\fancyhf{}
\rhead{\textcolor{secondary}{\textbf{ShieldIPTV} -- Guide UX / QA}}
\lhead{\textcolor{secondary}{webOS TV Submission}}
\rfoot{\thepage}
\renewcommand{\headrulewidth}{1pt}
\renewcommand{\headrule}{\hbox to\headwidth{\color{primary}\leaders\hrule height \headrulewidth\hfill}}

% Section title styling
\titleformat{\section}
  {\color{secondary}\normalfont\Large\bfseries}
  {\thesection}{1em}{}
\titlespacing*{\section}{0pt}{15pt}{10pt}

% Custom command for Scenarios
\newenvironment{scenario}[2]{%
  \begin{tcolorbox}[
    colback=light,
    colframe=primary,
    coltitle=white,
    title={\textbf{#1} : #2},
    fonttitle=\large\bfseries,
    arc=4px,
    boxrule=1.5pt,
    top=10pt,
    bottom=10pt,
    left=10pt,
    right=10pt
  ]
}{%
  \end{tcolorbox}
  \vspace{10pt}
}

\begin{document}

% Title / Header Layout
\begin{center}
    \vspace*{10pt}
    {\Huge \textbf{Shield\textcolor{primary}{IPTV}}} \\[5pt]
    {\Large \textbf{Guide des Scénarios d'Utilisation \& QA}} \\[5pt]
    {\large Soumission LG Content Store (webOS TV)}
    \vspace{20pt}
\end{center}

% Metadata Box
\begin{tcolorbox}[colback=light, colframe=border, arc=6px, boxrule=1pt]
\begin{tabular}{ll}
    \textbf{Nom de l'application :} & Shield IPTV Player \\
    \textbf{App ID :} & \texttt{com.shieldiptv.app} \\
    \textbf{Version de l'application :} & 1.0.8 (ou 1.0.0) \\
    \textbf{Mode de contrôle :} & Télécommande Magic Remote \& Touches directionnelles D-Pad
\end{tabular}
\end{tcolorbox}

\vspace{15pt}

\section{Présentation}
ShieldIPTV est une application client conçue pour le chargement et la lecture de listes de lecture IPTV appartenant à l'utilisateur (Xtream Codes API ou format de lien M3U).
\begin{itemize}[label=\textbullet]
    \item L'application \textbf{ne fournit et n'intègre aucun contenu multimédia}.
    \item Une \textbf{"Playlist de Démonstration"} intégrée est mise à disposition pour permettre aux testeurs QA de valider immédiatement les fonctionnalités de l'application avec des flux légaux du domaine public.
\end{itemize}

\vspace{15pt}

\section{Scénarios de test étape par étape}

\begin{scenario}{Scénario 1}{Lancement de l'application \& CGU}
\textbf{Action :} Lancer l'application ShieldIPTV sur le téléviseur ou l'émulateur webOS.\\
\textbf{Comportement attendu :}
\begin{itemize}[leftmargin=15pt, label=\textbullet]
    \item Un écran de démarrage s'affiche, puis la fenêtre des Conditions Générales d'Utilisation (CGU) s'ouvre.
    \item Le focus de la télécommande est automatiquement positionné sur le bouton \textbf{"Accepter" (Accepter)}.
    \item Le testeur peut modifier la langue de l'interface en utilisant le bouton de changement de langue (Français, Anglais, Espagnol, Italien).
    \item L'activation du bouton \textbf{"Accepter"} ferme la fenêtre et redirige vers le gestionnaire de listes de lecture.
\end{itemize}
\end{scenario}

\begin{scenario}{Scénario 2}{Connexion avec la Playlist Démo}
\textbf{Action :} Sur l'écran du gestionnaire de listes de lecture, naviguer vers le bas et sélectionner le bouton "Playlist Démo".\\
\textbf{Comportement attendu :}
\begin{itemize}[leftmargin=15pt, label=\textbullet]
    \item Un indicateur de chargement ("Connexion...") s'affiche à l'écran.
    \item Une fois la connexion établie, l'écran de transition vers le menu principal (\textbf{Portal}) s'affiche automatiquement.
\end{itemize}
\end{scenario}

\newpage

\begin{scenario}{Scénario 3}{Navigation sur le Portail principal}
\textbf{Action :} Navigation au sein du tableau de bord de l'application.\\
\textbf{Comportement attendu :}
\begin{itemize}[leftmargin=15pt, label=\textbullet]
    \item Trois grandes catégories sont proposées côte à côte : \textbf{LIVE TV}, \textbf{FILMS}, et \textbf{SÉRIES}.
    \item Des raccourcis d'outils sont présents en haut à droite : Gestion de compte, Test de débit, Vérificateur de lien, Test de flux, Paramètres, et Actualisation.
    \item L'utilisateur utilise les touches directionnelles \textbf{Gauche / Droite} pour basculer d'une catégorie à l'autre. Le focus est clairement matérialisé par une bordure lumineuse.
    \item Cliquer sur \textbf{LIVE TV} ouvre la grille des chaînes.
\end{itemize}
\end{scenario}

\begin{scenario}{Scénario 4}{Lecture en direct \& Plein écran}
\textbf{Action :} Sélection d'une catégorie et démarrage de la lecture d'une chaîne.\\
\textbf{Comportement attendu :}
\begin{itemize}[leftmargin=15pt, label=\textbullet]
    \item L'interface se sépare en une liste de catégories à gauche et une grille de chaînes à droite avec un mini-lecteur.
    \item Naviguer dans la grille et appuyer sur \textbf{OK} sur une chaîne. Le flux démarre dans le mini-lecteur de prévisualisation à droite.
    \item Appuyer de nouveau sur \textbf{OK} sur la chaîne sélectionnée (ou sur le mini-lecteur) pour activer le mode \textbf{Plein écran}.
    \item En plein écran, appuyer sur \textbf{OK} pour afficher les contrôles de lecture. Les touches \textbf{Haut / Bas} ouvrent la liste de zapping rapide.
    \item Appuyer sur la touche \textbf{RETOUR} (Back) de la télécommande pour fermer le plein écran et revenir à la grille.
\end{itemize}
\end{scenario}

\begin{scenario}{Scénario 5}{Test de débit réseau (Speed Test)}
\textbf{Action :} Naviguer sur le portail et lancer le test de débit.\\
\textbf{Comportement attendu :}
\begin{itemize}[leftmargin=15pt, label=\textbullet]
    \item Cliquer sur l'icône de test de débit (compteur de vitesse) dans la barre d'outils supérieure.
    \item Cliquer sur \textbf{"Lancer le test"} pour démarrer la mesure de bande passante descendante.
    \item La jauge s'anime et affiche le résultat final en Mbps.
    \item Appuyer sur la touche \textbf{RETOUR} pour revenir au Portail principal.
\end{itemize}
\end{scenario}

\vfill
\begin{center}
    \small\textcolor{gray}{Ce document sert de guide de validation d'expérience utilisateur (UX Scenario) pour la certification LG Content Store (webOS).}
\end{center}

\end{document}
